\documentclass[10pt, a4paper]{article}

\usepackage[left=1.4cm, top=1.4cm, right=1.4cm, bottom=1.4cm]{geometry}
\usepackage{amsmath,amssymb}
\usepackage{enumitem}
\usepackage{hyperref}
\usepackage{titlesec}
\usepackage{xcolor}
\usepackage{tabularx}
\usepackage{fancyhdr}
\usepackage{fontspec}
\usepackage{polyglossia}
\setmainlanguage{russian}
\setotherlanguage{english}
\setmainfont{Times New Roman}
\setsansfont{Arial}
\setmonofont{Consolas}

\definecolor{accent}{HTML}{333333}
\definecolor{graytext}{HTML}{666666}
\definecolor{linkblue}{HTML}{2563EB}

\hypersetup{colorlinks=true,linkcolor=linkblue,urlcolor=linkblue}

\titleformat{\section}{\large\bfseries\color{accent}\uppercase}{}{0em}{}[\vspace{-0.5em}\rule{\textwidth}{0.8pt}]
\titlespacing*{\section}{0pt}{8pt}{4pt}

\pagestyle{fancy}
\fancyhf{}
\renewcommand{\headrulewidth}{0pt}
\fancyfoot[R]{\small\color{graytext}\thepage}

\setlength{\parindent}{0pt}
\setlength{\parskip}{1pt}
\setlist[itemize]{nosep,leftmargin=1.2em,topsep=1pt}

\newcommand{\entry}[4]{{\textbf{#1}} \hfill {\small\color{graytext}#2}\\{\small #3 \hfill #4}\vspace{2pt}}
\newcommand{\projentry}[3]{{\textbf{#1}} {\small\color{graytext}| #2} \hfill {\small #3}\\}

\begin{document}

\begin{center}
{\LARGE\bfseries Артур Донской}\\[2pt]
{\normalsize Systems Engineer $\cdot$ Quantitative Engineer $\cdot$ 3D Engine Developer}\\[3pt]
{\small
\href{https://github.com/KorsarOfficial}{github.com/KorsarOfficial} $\cdot$
\href{https://leetcode.com/u/korsarofficial}{leetcode.com/u/korsarofficial} $\cdot$
Астрахань, Россия
}
\end{center}

\vspace{-2pt}

\section{Навыки}

\begin{tabularx}{\textwidth}{@{}r|X@{}}
\textbf{\small Языки} & C, C++20, Go, x86 Assembly, Python, JavaScript, ABAP, C\#, MATLAB \\
\textbf{\small Графика} & Vulkan API, OpenGL, SDL2, GLFW, Compute Shaders, FMod \\
\textbf{\small Системное} & AVX2 SIMD, lock-free SPSC, \texttt{alignas(64)}, acquire/release ordering \\
\textbf{\small Медиа} & FFmpeg (libavformat, libavcodec, libswscale), LZW, GIF89a \\
\textbf{\small Квант} & Байесовский вывод (Beta-Binomial), критерий Келли, Монте-Карло, EMA/ROC \\
\textbf{\small Инфра} & WebSocket, REST API, Telegram Bot API, SQLite, PostgreSQL, Docker, CMake, Git, Linux \\
\end{tabularx}

\vspace{-2pt}

\section{Опыт работы}

\entry{Quantitative Developer}{Фев 2026 -- н.в.}{Polyalphia}{Удалённо}
\begin{itemize}
\item Спроектировал торговую систему на Go (20 372 строки, 131 файл, 19 фаз разработки).
\item 5 стратегий: Mean Reversion (Sharpe 0.438, WR 53\%, \$112K PnL), Momentum (ROC+EMA, $\kappa$=0.25), Probability Fade (арбитраж, velocity filter), Time Decay (WR 85\%+), Longshot Fade (Griffith 1949).
\item Байесовский фильтр (Beta-Binomial, exp forgetting), Kelly sizing ($f_{\max}$=0.20), гистерезисный детектор.
\item Оркестратор, 2-уровневый риск (15\%/25\%, 40\% drawdown halt), WebSocket, сканер рынков, Telegram-бот, веб-дашборд, SQLite.
\end{itemize}

\entry{Системный программист (самозанятый)}{Окт 2023 -- н.в.}{Независимый}{Удалённо}
\begin{itemize}
\item KorsEngine: 3D Vulkan-движок на чистом C (1381 коммит, 148 файлов, 6.4 МБ). PBR, compute shaders, ручное управление памятью.
\item MP4-to-GIF (C++20): 4.3$\times$ ускорение (12.0с$\to$2.8с). Квантизация Ву (3D prefix sums), 32K AVX2 кэш (118М вызовов устранено), Floyd-Steinberg (zero bounds checks), собственный LZW (хэш-таблица 8192), lock-free SPSC (alignas(64)), single-decode ($-$56\%). Итог: 12\% от теоретического оптимума.
\item VideoCutter (C++/FFmpeg/SDL2), Korscord (C++, K-Means++), Snake Strategy (C++, A* на торе $\mathbb{Z}/40\mathbb{Z} \times \mathbb{Z}/30\mathbb{Z}$).
\end{itemize}

\entry{Разработчик ПО (самозанятый)}{Июл 2022 -- Окт 2023}{Независимый}{Удалённо}
\begin{itemize}
\item Управление роботом RM001-M03 (Python/PyQt5, USB/CH341): траектории, сканирование устройств, эмуляция.
\item GPU-физика на Taichi Lang. Прототипы на Unity (FPS, AR, RTS).
\end{itemize}

\vspace{-2pt}

\section{Open Source (16 PR)}

\projentry{Yandex Perforator}{Go, кластерный профайлер}{7 PR}
\begin{itemize}
\item $O(n){\to}O(\text{miss})$ eBPF syscalls, устранение гонки данных (RLock), exponential backoff, pre-size аллокаций, освобождение мьютекса во время Parse.
\end{itemize}

\projentry{Yandex Yatool}{Python/C++, билд-система}{5 PR}
\begin{itemize}
\item Старт $-$38мс, CPP\_FAKEID blast radius 82.7\%$\to$35\%, инстанцирование шаблонов $-$118с ($-$64.9\%), планирование $-$38.7с ($-$11.8\%), pre-size хэш-карты $-$0.26с.
\end{itemize}

\projentry{VK KPHP}{PHP/C++, компилятор}{4 PR}
\begin{itemize}
\item \texttt{instanceof} fix, \texttt{\_\_clone} parent chain, classmap autoloading, callback validation.
\end{itemize}

\vspace{-2pt}

\section{Образование}

\entry{Прикладная математика и информатика}{2025 -- 2029}{Московский физико-технический институт (МФТИ)}{Дистанционно}
{\small Кафедра системного программирования (ИСП РАН). Бакалавриат.}

{\small Самообразование: Кнут, <<Искусство программирования>>, тома 1--4A.}

\vspace{-2pt}

\section{Соревнования}

\begin{itemize}
\item \textbf{Smolathon 2025} --- тимлид, 19/80+ команд. Спроектировал архитектуру веб-платформы для ЦОДД Смоленской области. Прототип за 24ч очного финала. Декомпозиция задач, код-ревью, финальный питчинг. Форум <<Свой Код>>, призовой фонд 500К руб.
\item \textbf{Космический кубок: Миссия <<ЛУНА>>} --- планировщик лунной базы: анализ рельефа, размещение инфраструктуры с проверкой ограничений, мониторинг ресурсов. JS + WebGL визуализация, Python backend.
\item \textbf{LeetCode} --- 70 задач (22E, 34M, 14H). Общий acceptance rate 70.7\%. Hard acceptance \textbf{82.4\%} (14/17 решений принято с первой-второй попытки). Фокус: графы, DP, деревья, sliding window. \href{https://miro.com/app/board/uXjVGcc9d58=/?share_link_id=838442415695}{Доска декомпозиции паттернов на Miro}.
\end{itemize}

\vspace{-2pt}

\section{Языки}

\begin{tabularx}{\textwidth}{@{}r|X@{}}
\textbf{\small Русский} & Родной \\
\textbf{\small Английский} & Техническая документация, код, коммиты, PR-описания на английском. Чтение: свободно. \\
\end{tabularx}

\vspace{-2pt}

\section{Контакты}

\begin{tabularx}{\textwidth}{@{}r|X@{}}
\textbf{\small Telegram} & @korsarofficial \\
\textbf{\small GitHub} & \href{https://github.com/KorsarOfficial}{github.com/KorsarOfficial} \\
\textbf{\small LeetCode} & \href{https://leetcode.com/u/korsarofficial}{leetcode.com/u/korsarofficial} \\
\end{tabularx}

\end{document}
